\section{Introduction}

Diffusion transformers generate a video by iterating a denoiser fifty times, and classifier-free guidance doubles each iteration into a conditional and an unconditional pass, so a single clip costs a hundred forward passes. Training-free caching reduces this cost without touching the weights: a small number of steps are evaluated exactly, and the remaining steps reuse the last exact output. TeaCache, TaylorSeer and RA-CFGCache all follow this recipe and report two to three times fewer calls at a modest loss of quality.

The three methods differ in how a skipped step is reconstructed, and each decides when to skip from a quantity it has already observed. TeaCache accumulates the change of the modulated input and skips while that change stays under a level, TaylorSeer caches on a fixed interval and extrapolates in feature space, and RA-CFGCache composes a risk from the two guidance branches. The decision is therefore local in time, since it compares a per-step quantity against a fixed level, and the level is the same wherever the step sits in the trajectory.

We asked what a skipped step actually costs, and measured it by removing one step at a time. On CogVideoX-2B at thirty steps the cost is positional and steep. Skipping the first step costs 0.66 LPIPS against the full run, the second 0.41, the sixth 0.05, and each step after the tenth about 0.009, so the first ten steps carry ninety-two percent of the total. Two prompts rank the first eight steps identically and correlate at 0.94, which says the profile belongs to the sampler. An error introduced while the latent is still noise is carried by every later step, and an error introduced near the end is re-explained before the decode.

The profile turns placement into an allocation. With a budget of $B$ calls, the rule we use keeps the head exact and then grows the gaps geometrically along the tail, ending on the last step, which is the only step whose output is decoded directly. The length of the head follows from the measured profile instead of a search, and the number of calls equals the budget by construction, so the schedule is comparable with any other policy at the same cost. No part of it is learned and no calibration pass is needed beyond the one measurement.

We evaluate inside a single released pipeline, so that every arm shares the reference and the skip semantics, and our schedule enters through one hook that replaces the decision of whether to evaluate. At a matched number of calls our schedule reaches 0.068 LPIPS on CogVideoX-2B at seventeen frames where RA-CFGCache reaches 0.082 and its threshold-based cache 0.216, 0.037 against 0.043 at forty-nine frames, and 0.063 against 0.086 on Wan2.1. The gains come from placement alone: keeping the calls and replacing the schedule by a threshold loses a factor of three, while keeping the schedule and replacing the reconstruction by plain reuse changes nothing.
